\section{Quantum evolution and worst-case fidelity}\label{sec:quantum}

We verify directly that the two sector histories analyze
\eqref{eq:generator}. This also checks the continuous-time CSS
hypotheses in the forward memory implication of
Drechsler--Oreg--Aharonov, Theorem 1 (Section 1.5.1; formal version Section 3.1) \cite{DOA}, without requiring that
implication as a probabilistic input.

\begin{lemma}[Pauli-trajectory realization]\label{lem:trajectories}
For every finite \(L\), every encoded density operator \(\rho\),
possibly entangled with an arbitrary reference, and every \(t\ge0\),
the state \((e^{t\Lind_L}\otimes\id)(\rho)\) is the mixture over the
independent marked histories specified in
Section~\ref{sec:model}, applying their cumulative physical Paulis
to \(\rho\). Each sector has the single-face classical correction
rule of Section~\ref{sec:local} and noise rate at most \(2\eta\).
The probabilities of these histories are independent of \(\rho\).
\end{lemma}
\begin{proof}
On a Pauli translate \(U\rho U^\dagger\) of any encoded state, every
stabilizer has a definite sign determined by \(U\), regardless of
the logical state or the reference. Thus any tested projector
\(\Pi\) has a definite scalar value \(b\in\{0,1\}\) on this whole
branch. For \(J=P\Pi\), \(J^\dagger J=\Pi\), and
\[
 \Diss[P\Pi](U\rho U^\dagger)=
 \begin{cases}
 0,&b=0,\\
 P U\rho U^\dagger P-U\rho U^\dagger,&b=1.
 \end{cases}
\]
A rate-one attempt clock gives exactly this infinitesimal transition
law. A noise clock with rate \(\gamma_{p,P}\) gives
\(\gamma_{p,P}(P_p\rho P_p-\rho)\).
Summing these transition laws over the finitely many clocks shows
that the marked-history mixture solves the Lindblad equation with
the prescribed initial state. Finite-dimensional uniqueness of
this linear differential equation identifies it with the quantum
evolution. This uniqueness follows, for example, by the convergent
matrix exponential of the finite-dimensional superoperator.

Products of physical Paulis stay physical Paulis up to a phase.
Their syndrome signs depend only on the error chains, so neither
the correction rates nor the history probabilities depend on the
logical state. A primal correction adds its plaquette to the
primal error chain, and a dual correction adds its dual plaquette
to the dual error chain. The noise marginals are the sums of the
\(Z,Y\) and \(X,Y\) rates respectively. Shared \(Y\) marks can
correlate the sectors, but do not alter either marginal law.
\end{proof}

The stipulated generator is a correction Lindbladian in
\cite[Definition 2.24]{DOA}: its Hamiltonian vanishes, its \(Z\)
jumps use only \(X\)-check functions, and its \(X\) jumps use only
\(Z\)-check functions. Its noise has Definition 2.25's independent
Pauli form. The initial encoding has all check signs positive,
and \(D\) is a compatible decoder with \(D(0)=0\) as in their
Definition 2.5. These checks explain the applicability of their
forward Theorem 1 (Section 1.5.1; formal version Section 3.1). The direct realization just proved also handles
the stronger reference statement used here.

To implement \(\readout\), measure the commuting checks, apply
\(Z^{D_{\rm primal}(s_X)}\) and
\(X^{D_{\rm dual}(s_Z)}\), invert the code encoding on the corrected
code space, and output the designated qubit. The order of the two
Pauli corrections changes only a phase. Each correction removes
its associated syndrome and commutes with the other family of
checks, so their combination returns every attainable branch to
the code space.

\begin{proof}[Proof of Theorem~\ref{thm:main}]
Apply Corollary~\ref{cor:retention} to the two marginal sectors in
Lemma~\ref{lem:trajectories}. A union bound, requiring no independence
between sectors, gives an event \(\mathcal S_t\) with
\[
 \Prb(\mathcal S_t)\ge1-2C_0(1+t)e^{-\beta L}
\]
on which both filled classes remain zero through \(t\).
On such a history the primal and dual residual error chains are
boundaries of 3-chains. Their Paulis are products of \(Z\) and \(X\)
stabilizers, respectively. The total residual acts as the identity
on every code vector, up to a physically irrelevant phase.
It therefore restores all six logical qubits exactly, including
their joint state with a reference.

Since the history probabilities do not depend on the input,
the effective channel is a convex mixture whose contribution from
\(\mathcal S_t\) is exactly
\(\Prb(\mathcal S_t)\id\). For any pure \(|\Psi\rangle_{QR}\) the
target overlap on the other branches is nonnegative. Hence
\begin{equation}\label{eq:full-fidelity}
 \langle\Psi|(\logical\otimes\id_R)
  (|\Psi\rangle\langle\Psi|)|\Psi\rangle
 \ge1-2C_0(1+t)e^{-\beta L}.
\end{equation}
This conclusion is uniform in the fixed unused logical state
\(\sigma\); on success that state too is restored before it is
discarded.

Choose \(c=\beta/3\). For \(0\le t\le e^{cL}\),
\[
 2C_0(1+t)e^{-\beta L}
 \le4C_0e^{-(\beta-c)L}
 \le4C_0e^{-cL}.
\]
Take \(C=4C_0\), enlarging it if needed for a common \(L_0\).
With a trivial reference, minimizing \eqref{eq:full-fidelity} over
all pure inputs yields \eqref{eq:main}. With a general reference
it gives the asserted stronger inequality. All geometric and
probabilistic constants were fixed independently of \(L\), the
rate values within the bound, and the input. The fillings use only
the syndrome and fixed cell orderings. This proves every claimed
uniformity.
\end{proof}

For a Bell input, the same argument gives entanglement fidelity
at least \(1-Ce^{-cL}\). The maximum overlap of a pure two-qubit
product with a Bell vector is \(1/2\), by the Cauchy--Schwarz
inequality applied to its two Schmidt terms. Convexity extends
this upper bound to separable states. Thus for sufficiently
large \(L\) the guaranteed Bell fidelity is above that classical
benchmark throughout the interval. The worst-case statement in
Theorem~\ref{thm:main} does not rely on choosing this particular
input.
